\chapter{Bölüm Topolojisi}

Bölüm topolojisi, bir topolojik uzayın noktalarını denklik sınıflarına göre
\emph{yapıştırarak} yeni bir uzay kurar.
Alt uzay ve çarpım topolojisiyle birlikte temel üç inşaat yönteminden birini oluşturur.

%-------------------------------------------------------------------
\section{Konu}

\begin{definition}[Bölüm Kümesi]
$(X, \tau)$ topolojik uzayı ve $X$ üzerinde bir denklik bağıntısı $\sim$ verilsin.
Denklik sınıfı $[x] = \{y \in X : y \sim x\}$.
\textbf{Bölüm kümesi:} $X/{\sim} = \{[x] : x \in X\}$.
\end{definition}

\begin{definition}[Bölüm Topolojisi]
Bölüm haritası $q: X \to X/{\sim}$,\ $q(x) = [x]$.
\textbf{Bölüm topolojisi} $\tau_{X/\sim}$ üzerinde:
\[
U \subseteq X/{\sim} \text{ açıktır} \iff q^{-1}(U) \in \tau.
\]
Bu, $q$'yu sürekli kılan en ince topolojidir.
\end{definition}

\begin{sezgi}
Bölüm topolojisini bir ``katlama'' gibi düşünün. Elinizde bir kâğıt şerit
($[0,1]$) varsa ve iki ucunu birbirine yapıştırırsanız bir çember ($S^1$) elde
edersiniz. Yapıştırma kuralı bir denklik bağıntısıdır ($0 \sim 1$); bölüm
topolojisi ise ``kâğıdı yırtmadan'' yapılan bu katlamanın doğal topolojisidir.
$q$ haritası her noktayı yapıştırma sonrası bulunduğu konuma gönderir; bir
kümenin yapıştırılmış uzayda açık olması, ön-görüntüsünün orijinal kâğıtta açık
olmasına bağlıdır.
\end{sezgi}

\begin{definition}[Doymuş Küme]
Bir $A \subseteq X$ kümesi \textbf{doymuştur} (saturated), eğer her $x \in A$
için $x$'in tüm denklik sınıfı $[x]$ yine $A$ içinde kalıyorsa, yani
$A = q^{-1}(q(A))$ ise. Bölüm topolojisinin açık kümeleri tam olarak $X$'in
doymuş açık kümelerine karşılık gelir:
\[
q(A)\ X/{\sim}\ \text{içinde açık} \iff A\ \text{doymuş ve}\ X\ \text{içinde açık.}
\]
\end{definition}

\begin{figure}[ht]
\centering
\begin{tikzpicture}[font=\small, etiket/.style={font=\scriptsize}]
  % SOL: [0,1] serit -- uclar 0 ve 1 isaretli
  \draw[blue!70!black, very thick] (0,0) -- (4,0);
  \fill[red!75!black] (0,0) circle (2.2pt);
  \fill[red!75!black] (4,0) circle (2.2pt);
  \node[red!70!black, etiket, below=3pt] at (0,0) {$0$};
  \node[red!70!black, etiket, below=3pt] at (4,0) {$1$};
  \node[etiket, above=4pt] at (2,0) {$[0,1]$ şerit};

  % yapistirma oku
  \draw[-{Stealth[length=2.6mm]}, orange!85!black, very thick]
    (4.6,0) -- (6.2,0);
  \node[orange!80!black, etiket, align=center, above=2pt] at (5.4,0.1)
    {$0 \sim 1$\\yapıştır};

  % SAG: cember S^1, yapistirma noktasi tek
  \draw[blue!70!black, very thick] (8.4,0) circle (1.1);
  \fill[red!75!black] (8.4,1.1) circle (2.4pt);
  \node[red!70!black, etiket, above=3pt] at (8.4,1.1) {$[0]=[1]$};
  \node[etiket, below=4pt] at (8.4,-1.15) {$S^1 = [0,1]/{\sim}$};

  % q haritasi etiketi (cemberin altindaki etikete carpmasin diye sola hizali)
  \node[etiket, align=center, fill=blue!5, rounded corners, inner sep=4pt,
    text width=6.2cm, anchor=north west] at (-0.2,-2.3)
    {Bölüm haritası $q$ her noktayı sınıfına gönderir; \emph{tek} fark, iki ucun
     tek bir noktada (\,$[0]=[1]$\,) çakışmasıdır.};
\end{tikzpicture}

\caption{$[0,1]$ şeridinin uçları yapıştırılarak çember $S^1$ elde edilir;
$q$ haritası her noktayı denklik sınıfına gönderir.}
\end{figure}

\begin{table}[h]
\centering
\begin{tabular}{ll}
\toprule
\textbf{Özellik} & \textbf{Bölüm uzayı için durum} \\
\midrule
Kompaktlık & $X$ kompakt $\Rightarrow X/{\sim}$ kompakt \\
Bağlantılılık & $X$ bağlantılı $\Rightarrow X/{\sim}$ bağlantılı \\
Hausdorff & Genel olarak korunmaz \\
T1 & $\sim$ kapalı bağıntı ise korunur \\
\bottomrule
\end{tabular}
\caption{Topolojik özelliklerin bölüm altında korunumu}
\end{table}

\begin{figure}[ht]
\centering
\begin{tikzpicture}[font=\small, etiket/.style={font=\scriptsize},
  nokta/.style={circle, fill=blue!70!black, inner sep=1.8pt}]
  % SOL: X = {a,b,c,d}, sinif kutulariyla gruplanmis
  \node[etiket, above] at (1.3,2.2) {$X = \{a,b,c,d\}$};

  % {a,b} kutusu
  \draw[violet!60!black, thick, rounded corners, fill=violet!7]
    (-0.1,1.0) rectangle (1.2,2.0);
  \node[nokta, label={[etiket]left:$a$}] (a) at (0.25,1.5) {};
  \node[nokta, label={[etiket]right:$b$}] (b) at (0.85,1.5) {};
  \node[violet!60!black, etiket, below=1pt] at (0.55,1.0) {$\{a,b\}$};

  % {c}
  \draw[teal!55!black, thick, rounded corners, fill=teal!7]
    (1.7,1.15) rectangle (2.5,1.85);
  \node[nokta, label={[etiket]above:$c$}] (c) at (2.1,1.5) {};

  % {d}
  \draw[orange!75!black, thick, rounded corners, fill=orange!7]
    (3.0,1.15) rectangle (3.8,1.85);
  \node[nokta, label={[etiket]above:$d$}] (d) at (3.4,1.5) {};

  % q oklari asagi
  \draw[-{Stealth[length=2.2mm]}, gray!70] (0.55,0.95) -- (0.55,-0.05);
  \draw[-{Stealth[length=2.2mm]}, gray!70] (2.1,1.1) -- (2.1,-0.05);
  \draw[-{Stealth[length=2.2mm]}, gray!70] (3.4,1.1) -- (3.4,-0.05);
  \node[etiket, right] at (3.55,0.55) {$q$};

  % SAG/ALT: X/~ = uc nokta
  \node[etiket, below] at (1.95,-0.9) {$X/{\sim} = \{[a],[c],[d]\}$ (üç nokta)};
  \node[circle, draw=violet!60!black, fill=violet!18, inner sep=2.6pt,
    label={[etiket]below:$[a]$}] at (0.55,-0.45) {};
  \node[circle, draw=teal!55!black, fill=teal!18, inner sep=2.6pt,
    label={[etiket]below:$[c]$}] at (2.1,-0.45) {};
  \node[circle, draw=orange!75!black, fill=orange!18, inner sep=2.6pt,
    label={[etiket]below:$[d]$}] at (3.4,-0.45) {};
\end{tikzpicture}

\caption{Denklik sınıfları bölüm uzayında birer noktaya çöker: dört nokta
$\{a,b\}, \{c\}, \{d\}$ sınıflarıyla üç noktaya iner.}
\end{figure}

\begin{karsiornek}
Hausdorff özelliği bölüm altında korunmaz. $\mathbb{R}$ üzerinde
``$x \sim y \iff x - y \in \mathbb{Q}$'' bağıntısını alın. Bölüm uzayı
$\mathbb{R}/\mathbb{Q}$ kaba (indiscrete) bir uzaydır: tek açık kümeler
$\emptyset$ ve tüm uzaydır, çünkü $\mathbb{Q}$ ile öteleme altında doymuş tek
açık küme boş küme ve $\mathbb{R}$'dir. Dolayısıyla iki farklı nokta hiçbir
zaman ayrık komşuluklarla ayrılamaz --- $\mathbb{R}$ Hausdorff olmasına rağmen
$\mathbb{R}/\mathbb{Q}$ T1 bile değildir. Sonlu modelde de aynı tuzak geçerlidir:
doymamış açık kümeler yapıştırma sonrası ``kaybolur'' ve ayırma aksiyomlarını
bozabilir.
\end{karsiornek}

%-------------------------------------------------------------------
\section{Teoremler}

\begin{theorem}[Bölüm haritası sürekliliği]
$q: X \to X/{\sim}$ her zaman süreklidir.
\end{theorem}

\begin{theorem}[Evrensel Özellik]
$g: X/{\sim} \to Z$ olsun.
$g$ süreklidir $\Longleftrightarrow g \circ q: X \to Z$ süreklidir.
\begin{proof}[İspat eskizi]
($\Rightarrow$) $g$ sürekli ve $q$ sürekli (Teorem~1) olduğundan, iki sürekli
haritanın bileşkesi $g \circ q$ süreklidir.
($\Leftarrow$) $g \circ q$'nun sürekli olduğunu varsayalım. $Z$'de bir $V$ açık
kümesi alın; $g^{-1}(V) \subseteq X/{\sim}$ kümesinin bölüm topolojisinde açık
olduğunu göstermeliyiz. Tanım gereği bu, $q^{-1}(g^{-1}(V))$ kümesinin $X$'te
açık olmasına denktir. Fakat
$q^{-1}(g^{-1}(V)) = (g \circ q)^{-1}(V)$ ve $g \circ q$ sürekli olduğundan bu
küme $X$'te açıktır. Dolayısıyla $g^{-1}(V)$ açıktır ve $g$ süreklidir. Bu
özellik, bölüm uzayından çıkan sürekli haritaların kontrolünü $X$ üzerindeki
haritaların kontrolüne indirger ve bölüm topolojisini evrensel bir nesne olarak
karakterize eder.
\end{proof}
\end{theorem}

\begin{theorem}[Kompaktlık Korunumu]
$X$ kompakt $\Rightarrow X/{\sim}$ kompakttır.
\begin{proof}[İspat eskizi]
$q$ sürekli ve örtendir; kompakt bir kümenin sürekli görüntüsü kompakt olduğundan
$X/{\sim} = q(X)$ kompakttır.
\end{proof}
\end{theorem}

\begin{theorem}[Bağlantılılık Korunumu]
$X$ bağlantılı $\Rightarrow X/{\sim}$ bağlantılıdır.
\end{theorem}

\begin{theorem}[Hausdorff Kriteri]
$X/{\sim}$ Hausdorff $\Longleftrightarrow$ $\sim$ grafiği $X \times X$'te kapalıdır.
\end{theorem}

%-------------------------------------------------------------------
\section{Algoritmalar}

\subsection{quotient\_set — $O(|X|^2 \cdot \alpha(|X|))$}

\begin{algorithm}
\caption{QuotientSet$(X, \sim)$}
\begin{algorithmic}[1]
\State Her $x$ için sınıf $\leftarrow \{x\}$ (union-find başlat)
\For{each $(x, y) \in {\sim}$}
    \State \textbf{Union}$(x, y)$
\EndFor
\State \Return kök temsilcilere göre sınıfları tuple olarak
\end{algorithmic}
\end{algorithm}

\subsection{finite\_quotient\_contract — $O(|X| + |P|)$}

\begin{algorithm}
\caption{FiniteQuotientContract$(X, P)$}
\begin{algorithmic}[1]
\State \textit{// $P$: $X$'in örtüşmeyen bölüntüsü}
\State Doğrula: $P$ her $x$'i kapsar ve bloklar ayrık
\State \Return FiniteConstructionContract(status=\texttt{true}, carrier\_size=$|X|$, block\_count=$|P|$)
\end{algorithmic}
\end{algorithm}

%-------------------------------------------------------------------
\section{pytop API}

\begin{lstlisting}[language=Python]
from pytop import (
    quotient_set,               # denklik bagintisindan bolum kumesi
    finite_quotient_contract,   # boluntuden insaat sozlesmesi
    finite_quotient_summary,    # kisa ozet dizesi
    make_quotient_map,          # QuotientMap nesnesi
    is_quotient_map,            # Result: harita bolum haritasi mi?
    equivalence_class,          # bir noktanin denklik sinifi
    partition_from_equivalence, # denklik bagintisi -> boluntu
    is_equivalence_relation,    # baginti denklik bagintisi mi?
    equivalence_from_partition, # boluntu -> denklik bagintisi
    equivalence_from_classes,   # sinif bloklarindan denklik bagintisi
    canonical_projection_from_equivalence,  # q haritasini sozluk olarak
    discrete_topology,          # ayrik topoloji
    sierpinski_space,           # Sierpinski uzayi
    make_topology,              # tasiyici + aciklardan FiniteTopologicalSpace
)
\end{lstlisting}

\texttt{quotient\_set(carrier, relation)} $\to$ \texttt{tuple[frozenset, ...]}

\texttt{finite\_quotient\_contract(carrier, partition)} $\to$ \texttt{FiniteConstructionContract}

\texttt{finite\_quotient\_summary(carrier, partition)} $\to$ \texttt{str}

\texttt{make\_quotient\_map(domain, codomain)} $\to$ \texttt{QuotientMap}

\texttt{is\_quotient\_map(map\_obj)} $\to$ \texttt{Result}

\texttt{equivalence\_from\_partition(universe, partition)} $\to$ \texttt{set[tuple]}

\texttt{partition\_from\_equivalence(carrier, relation)} $\to$ \texttt{set[frozenset]}

\texttt{equivalence\_class(carrier, relation, point)} $\to$ \texttt{set}

\texttt{canonical\_projection\_from\_equivalence(carrier, relation)} $\to$ \texttt{dict}

%-------------------------------------------------------------------
\section{Örnekler}

\subsection*{Örnek 5.1 — İki Noktayı Özdeşleştirme}

\begin{lstlisting}[language=Python]
d4 = discrete_topology(0, 1, 2, 3)
partition_1 = [[0, 1], [2], [3]]
qs1 = quotient_set(
    [0, 1, 2, 3],
    [(0,0),(1,1),(2,2),(3,3),(0,1),(1,0)]
)
print("Bolum kumesi:", qs1)
contract1 = finite_quotient_contract([0, 1, 2, 3], partition_1)
print("Blok sayisi:", contract1.block_count)
print("Durum:", contract1.status)
\end{lstlisting}

\begin{verbatim}
Bolum kumesi: (frozenset({2}), frozenset({3}), frozenset({0, 1}))
Blok sayisi: 3
Durum: true
\end{verbatim}

\subsection*{Örnek 5.2 — Tüm Noktaları Tek Noktaya Çökertme}

\begin{lstlisting}[language=Python]
carrier = [0, 1, 2, 3]
qs2 = quotient_set(
    carrier,
    [(i, j) for i in carrier for j in carrier]
)
print("Tek blok:", qs2)
contract2 = finite_quotient_contract(carrier, [carrier])
print("Blok sayisi:", contract2.block_count)
r2 = contract2.to_result()
print("Metadata:", r2.metadata['block_sizes'])
\end{lstlisting}

\begin{verbatim}
Tek blok: (frozenset({0, 1, 2, 3}),)
Blok sayisi: 1
Metadata: [4]
\end{verbatim}

\subsection*{Örnek 5.3 — Özel Topolojide Bölüm Sözleşmesi}

\begin{lstlisting}[language=Python]
X5 = [1, 2, 3, 4, 5]
tau5 = [set(), {1,2}, {3,4}, {1,2,3,4}, {1,2,3,4,5}]
fts5 = make_topology(X5, *tau5)
partition5 = [[1, 2], [3, 4], [5]]
contract5 = finite_quotient_contract(X5, partition5)
r5 = contract5.to_result()
print("Status:", r5.status)
print("Bloklar:", r5.metadata['block_count'])
print("Blok boyutlari:", r5.metadata['block_sizes'])
\end{lstlisting}

\begin{verbatim}
Status: true
Bloklar: 3
Blok boyutlari: [2, 2, 1]
\end{verbatim}

\subsection*{Örnek 5.4 — Bölüm Haritası Nesnesi}

\begin{lstlisting}[language=Python]
d3 = discrete_topology(0, 1, 2)
d2 = discrete_topology('a', 'b')
qmap = make_quotient_map(d3, d2)
print("Etiketler:", sorted(qmap.tags))
r_qmap = is_quotient_map(qmap)
print("is_quotient_map status:", r_qmap.status)
print("Justification:", r_qmap.justification[0])
\end{lstlisting}

\begin{verbatim}
Etiketler: ['continuous', 'quotient', 'surjective']
is_quotient_map status: true
Justification: Map tag 'quotient' is explicitly present.
\end{verbatim}

\subsection*{Örnek 5.5 — [0,1] Uçlarını Yapıştırarak Çember $S^1$}

$[0,1]$ aralığını $5$ noktayla modelliyoruz. İki ucu ($0$ ve $4$)
özdeşleştirmek, şeridi bir çembere yapıştırmaya karşılık gelir.

\begin{lstlisting}[language=Python]
ring = [0, 1, 2, 3, 4]
ring_classes = [[0, 4], [1], [2], [3]]
ring_rel = equivalence_from_partition(ring, ring_classes)
print("Denklik mi?", is_equivalence_relation(ring, ring_rel))
print("Bolum kumesi boyutu:", len(quotient_set(ring, ring_rel)))
print("Ozet:", finite_quotient_summary(ring, ring_classes))
print("0'in sinifi:", sorted(equivalence_class(ring, ring_rel, 0)))
print("4'un sinifi:", sorted(equivalence_class(ring, ring_rel, 4)))
\end{lstlisting}

\begin{verbatim}
Denklik mi? True
Bolum kumesi boyutu: 4
Ozet: quotient: status=true, carrier=5, blocks=4
0'in sinifi: [0, 4]
4'un sinifi: [0, 4]
\end{verbatim}

Uçlar tek bir sınıfta birleşti: $5$ noktalı şerit, $4$ noktalı çembere indi.

\subsection*{Örnek 5.6 — Doymuş Kümeler ve Kanonik İzdüşüm}

Kanonik izdüşüm $q$ her noktayı denklik sınıfına gönderir. Bir kümenin
\textbf{doymuş} olması, içerdiği her noktanın sınıfını da tamamen içermesi
demektir.

\begin{lstlisting}[language=Python]
sat_carrier = ['a', 'b', 'c', 'd']
sat_part = [['a', 'b'], ['c'], ['d']]
sat_rel = equivalence_from_partition(sat_carrier, sat_part)
proj = canonical_projection_from_equivalence(sat_carrier, sat_rel)
for pt in sat_carrier:
    print(f"q({pt}) =", sorted(proj[pt]))
cls_a = equivalence_class(sat_carrier, sat_rel, 'a')
A = {'a', 'b'}
print("{a,b} doymus mu?",
      all(equivalence_class(sat_carrier, sat_rel, p) <= A for p in A))
print("{a} doymus mu?", cls_a <= {'a'})
\end{lstlisting}

\begin{verbatim}
q(a) = ['a', 'b']
q(b) = ['a', 'b']
q(c) = ['c']
q(d) = ['d']
{a,b} doymus mu? True
{a} doymus mu? False
\end{verbatim}

$\{a, b\}$ doymuştur, ama $\{a\}$ doymuş değildir: sınıf arkadaşı $b$'yi
içermez. Bölüm topolojisinin açık kümeleri ancak doymuş açık kümelerden gelir.

\subsection*{Örnek 5.7 — Sınıflardan Noktalara: Round-Trip}

\texttt{equivalence\_from\_classes} sınıf bloklarından denklik bağıntısı kurar;
\texttt{partition\_from\_equivalence} bağıntıyı tekrar bölüntüye çevirir.

\begin{lstlisting}[language=Python]
rt_blocks = equivalence_from_classes([1, 2], [3], [4, 5])
print("Denklik mi?", is_equivalence_relation([1, 2, 3, 4, 5], rt_blocks))
rt_part = partition_from_equivalence([1, 2, 3, 4, 5], rt_blocks)
print("Blok sayisi:", len(rt_part))
print("Bloklar:", sorted(sorted(b) for b in rt_part))
\end{lstlisting}

\begin{verbatim}
Denklik mi? True
Blok sayisi: 3
Bloklar: [[1, 2], [3], [4, 5]]
\end{verbatim}

Sınıflardan kurulan bağıntı, bölüntüye çevrilip geri okunduğunda aynı üç bloğu
verir: bölüm uzayının $3$ noktası, tam olarak $3$ denklik sınıfına karşılık gelir.

%-------------------------------------------------------------------
\section{Alıştırmalar}

\subsection*{Kodlama}

\begin{enumerate}[label=K\arabic*.]
\item \texttt{discrete\_topology(0,1,2,3,4)} üzerinde \texttt{[[0,4],[1,3],[2]]}
      bölüntüsünü kullanarak \texttt{finite\_quotient\_contract} çağırın;
      blok sayısını ve boyutlarını yazdırın.
\item \texttt{sierpinski\_space()} üzerinde trivial bölüntü ile
      \texttt{finite\_quotient\_summary} çağırın ve çıktıyı yorumlayın.
\item \texttt{make\_topology([1,2,3,4], set(), \{1,2\}, \{3,4\}, \{1,2,3,4\})} üzerinde
      \texttt{[[1,2],[3,4]]} ve \texttt{[[1],[2],[3],[4]]} bölüntülerini karşılaştırın.
\item \texttt{[0,1,2,3,4,5]} taşıyıcısı üzerinde \texttt{[[0,5],[1],[2],[3],[4]]}
      bölüntüsünü \texttt{equivalence\_from\_partition} ile denklik bağıntısına
      çevirin; \texttt{is\_equivalence\_relation} ile doğrulayın ve $0$ ile $5$'in
      \texttt{equivalence\_class} çıktılarının aynı olduğunu gösterin.
\item \texttt{['x','y','z']} taşıyıcısı ve \texttt{[['x','y'],['z']]} bölüntüsü için
      \texttt{canonical\_projection\_from\_equivalence} çağırıp her noktanın
      görüntüsünü yazdırın; \texttt{\{'x','y'\}} kümesinin doymuş,
      \texttt{\{'x'\}} kümesinin doymuş olmadığını gösterin.
\end{enumerate}

\subsection*{Teori}

\begin{enumerate}[label=T\arabic*.]
\item $q: X \to X/{\sim}$ bölüm haritasının $\tau_{X/\sim}$ tanımı gereği
      her zaman sürekli olduğunu ispatlayın.
\item $X$ bağlantılı $\Rightarrow X/{\sim}$ bağlantılı olduğunu ispatlayın.
      ($q$ sürekli ve surjektif; bağlantılı kümenin sürekli görüntüsü bağlantılıdır.)
\item Bölüm haritasının \textbf{evrensel özelliğini} ispatlayın:
      $g: X/{\sim} \to Z$ için $g$ sürekli $\iff g \circ q$ sürekli.
      \ipucu{$(g \circ q)^{-1}(V) = q^{-1}(g^{-1}(V))$ eşitliğini bölüm
      topolojisinin tanımıyla birleştirin.}
\item $A \subseteq X$ doymuş ($A = q^{-1}(q(A))$) ise $q(A)$'nın bölüm uzayında
      açık olması için $A$'nın $X$'te açık olmasının yeterli ve gerekli olduğunu
      gösterin. Doymamış bir $A$ için bu yazışmanın neden bozulduğunu açıklayın.
\end{enumerate}
